\section*{Résumé}
Nous étudions la partition d'une séquence en segments contigus par minimisation
exacte de l'inertie intra-classe. La méthode de Fisher est comparée aux K-means
et à la classification ascendante hiérarchique de Ward. Sur les données simulées,
le coude du critère conduit à quatre segments. Sur \file{aiguillage}, Fisher
retrouve les quatre classes de référence, avec un ARI de 1, contre $0{,}943$
pour les K-means à 20 départs et $0{,}859$ pour Ward. Les analyses
complémentaires montrent toutefois que la comparaison dépend de l'initialisation
des K-means et de l'échelle des variables.
